\section{Return of the previous origin}\label{sec:recovery}

We ask how the previous route returns after an invalid origin displaces a valid one.
For identifier gaps of at least 15 minutes, the return shows more of its final set at 30\,s than the outward change.
Its day-block intervals are above zero in 2022, 2023 and 2025.
But most matched pairs are flips under one minute apart.
In these flips the observed return set is larger than the outward set, and it shows less of its final set at 30\,s.

\begin{table*}[t]
\small
\centering
\caption{Matched outward ($O$) and return ($R$) transitions per year. Gap: time between the identifiers of the two events of a link. $\Delta\Jac$: paired difference of the 30-s \floor, $R$ minus $O$, with 95\,\% component-bootstrap intervals. These intervals treat links as almost independent, and the wider day-block intervals are in \Cref{tab:recovery-dayblock}. Return slower: share of links where $R$ has the larger $t_{90}$. ROV share: share of RoVista-measured AS--link observations that are RoVista-protected.}
\label{tab:recovery-summary}
\setlength{\tabcolsep}{3.5pt}
\begin{tabular}{lrrrrrrrrr}
\toprule
year & \hd{retained\\links} & \hd{links with\\gap $<1$\,min} & \hd{$\Delta\Jac$, all links\\{[95\,\% CI]}} & \hd{$\Delta\Jac$, gap $\geq 15$\,min\\{[95\,\% CI]}} & \hd{$t_{90}$ p50 (s)\\$O$ / $R$} & \hd{return\\slower} & \hd{extent\\$R/O$, p50} & \hd{$|R\setminus O|/|R|$\\mean} & \hd{ROV share\\$R\setminus O$ / $R\cap O$} \\
\midrule
\csname @@input\endcsname generated/recovery_summary_rows.tex % primitive input: expandable inside the tabular
\bottomrule
\end{tabular}
\end{table*}

% layout (round-3 reviews S6, second reviewer presentation): the pooled figure fig:pairs repeats Table 3 and moved unchanged to sections/app-tables.tex.
\paragraph{Matched transitions.}
Each outward transition ($O$) is a selected origin change in which an invalid origin displaces a valid one.
From 2023 on, it comes from the catalogue population.
In 2022, without catalogue labels, the outward changes come mostly from the label-free stratum, because the stratified sample holds only 96.
% src: docs/paper/corrected-recovery-2026-10-16.json years.2022.pairs.exclude_reason null and full_selection_links_all_populations = identified_links = 1,999;
%   years.2023--2026.pairs.exclude_reason "rpki_invalid_sample" (4,940 all-population links against 3,141 identified in 2023);
%   paper/generated/propagation-facts.json years.2022.stratified.invalid_displacing_valid 96, years.2022.invalid_stratum.invalid_displacing_valid 2,593;
%   recovery_population.tex 2022: 2,710 selected outward changes (round-3 review NEW-5).
We pair it with the return transition ($R$) that restores the previous origin within 24\,h (\Cref{sec:data}).
Each pair is a \emph{link}, and we measure both directions with the same rule.

For each collector, peer and prefix, we take the first observed transition from the preceding origin to the target origin at or after its identifier.
We keep the ASes on the paths of that transition, with a clock that starts at its first observation.
We do not count ordinary announcements of the valid origin before its return, so routes that never left do not inflate the return.

We exclude links with missing source files, no qualifying transition, under two hours of follow-up, or a return before its outward change.
Of \num{20464} identified links, \num{20256} remain.
% src: paper/generated/recovery_population.tex (identified 1,999 + 3,141 + 5,429 + 6,689 + 3,206 = 20,464; retained 1,982 + 3,109 + 5,355 + 6,684 + 3,126 = 20,256;
%   excluded 173 + 35 = 208). Corrected from 21,464, an addition error (round-3 review NEW-6 / X1).

A return event can belong to more than one link.
Intervals thus resample connected components of shared events, with 1,000 draws.
Almost every component holds one link, so a revision analysis also resamples UTC days of the outward identifier (\Cref{tab:recovery-dayblock}).
\Cref{fig:pairs} shows the pooled distributions.

\paragraph{All links.}
The mean paired difference in the 30-s \floor, return minus outward, is negative in every year (\Cref{tab:recovery-summary}).
It is $-0.05$ to $-0.16$ in 2023--2026, with intervals below zero also under day blocks, and $-0.005$ in 2022, with an interval that includes zero.
The return reaches 90\,\% of its two-hour extent later than the outward change in 53--63\,\% of links.
The median return $t_{90}$ is 56.5--88.1\,s.
The median ratio of return extent to outward extent is 1.06--1.54.
The signs stay the same when we exclude links whose source windows cross midnight, where the daily history of a session resets.
% src: docs/paper/corrected-recovery-2026-10-16.json years.*.same_day_sensitivity (-0.006 / -0.054 / -0.087 / -0.166 / -0.047; only 2022 includes zero)
% src: paper/generated/recovery_summary_rows.tex t90 column, R values 87.6 / 88.1 / 79.0 / 73.6 / 56.5 (round-2 review N15: report 56.5, not 56)

% layout: generated/recovery_gap_wide.tex is generated/recovery_gap.tex transposed by tools/recovery_gap_wide.py,
%   with every cell copied verbatim (the script checks the cell multisets); generated/recovery_gap.tex is unchanged.

\paragraph{By identifier gap.}
Flips with an identifier gap under one minute are 55--65\,\% of links (\Cref{tab:recovery-gap} in the appendix).
For them the paired difference is $-0.10$ to $-0.23$ with all intervals below zero, and the return is slower in 65--75\,\% of links.
For these flips, the median extent ratio is 1.09--1.57, and one UTC day holds 24\,\% of the 2025 flips (\Cref{tab:recovery-dayblock}).

Three reasons, which we did not test, can explain why the observed return set is larger.
First, the return can start before the outward route has finished spreading.
Second, ROV can keep the invalid outward route away from ASes that later carry the valid return.
Third, path exploration after the withdrawal of the invalid route can add transient paths to the return set, which keeps every AS once observed~\cite{labovitz2000delayed}.
Withdrawals also reach the collectors later than announcements (\Cref{sec:data}), which can delay the return.
% round 5 (review-r5 W3, item 2; panel-r5 W7(c), W11): the larger return set is a descriptive, secondary result; "We did not classify
%   what triggers each return, so these reasons are hypotheses" is folded into the first sentence ("which we did not test"), and the accumulation
%   caveat is explicit (sections/03-data.tex: "The set accumulates: an AS stays in it once observed").

Links with a gap of at least 15 minutes are 13--31\,\% of the total.
For them the difference is $+0.03$ to $+0.13$, and all component intervals are above zero.
Under day blocks, the intervals of 2024 and 2026 include zero.
In each of these years, one UTC day holds many such links with a negative mean (\Cref{tab:recovery-dayblock}).
A post hoc check, run for these two years only, gives a mean of $+0.09$ and $+0.10$ without that day.

The median return $t_{90}$ of these links is 30--56\,s, against 55--186\,s for the outward change.
Their median extent ratio is 0.95--1.26, closer to one than for the flips.
In 2023--2026 its point estimate stays above one.
The identifier gap is the time between two identified events, not a measured lifetime of the outward route.
These ratios therefore do not rule out an early end of the outward change.
% round 5 (review-r5 W3, item 2): the former clause "which an early end of the outward change alone does not explain" is removed; a long
%   identifier gap does not show a long continuous exposure of the outward route. The leave-largest-day-out values above are a post hoc check
%   of 2024 and 2026 only (panel-r5 W7(a), DA #13), not a robustness check of every stratum; their src is the comment below.
% src: paper/generated/recovery_gap.tex; gap strata: scripts/paper_recovery_dayblock.py GAPS (< 1 min = [0, 60) s; >= 15 min = [900, 86,400) s).
% src: paper/generated/recovery_dayblock.json years.2025.lt1min (links 4,306, max_links_per_unit 1,013); years.2024.gt15min and years.2026.gt15min
%   (links 1,259 / 809, max_links_per_unit 116 / 107, mean_diff.point +0.0336 / +0.0616, dayblock_excludes_zero false). Dates and means of the
%   largest day and the means without it: recovery_dayblock.json years.{2024,2026}.gt15min.largest_day (written by
%   scripts/paper_recovery_dayblock.py since round 7, review-r6 W5 "generate the two-year diagnostic"; raw UTC day hijack_id[:8]):
%   2024-10-17: 116 links, mean -0.5452, without it 1,143 links, mean +0.0924; 2026-08-01: 107 links, mean -0.1827, without it 702 links, mean +0.0989
%   (the same values as the earlier read-only recomputation; every other field of the JSON is unchanged).
% src: extent ratio >= 15 min, 2023--2026: 1.04 / 1.14 / 1.13 / 1.26 (recovery_gap.tex; recovery_dayblock.json median_ratio.point).
% src (round 4, panel-r4 NEW-12): extent ratio >= 15 min, day-block CI 2023 [1.00, 1.09], 2024 [1.05, 1.26], 2025 [1.07, 1.18],
%   2026 [1.08, 1.80] (recovery_dayblock.tex); Delta J >= 15 min day-block CI above zero in 2022, 2023, 2025 (same file), replacing "in most years".
% src: return hypotheses: round-3 reviews NEW-7 (truncation, ROV, path exploration); ROV: "Overlap and ROV" below.
% src: trigger classification of returns not done (docs/paper/camera-ready-20261030/reviews-20261030/RESPONSE-ROUND1.md "Not done")

\paragraph{Overlap and ROV\@.}
On average, 40--52\,\% of each return footprint is absent from its outward footprint.
Absence from observed paths does not show that an AS lacked or rejected a route.
The data therefore do not identify an outage in $R\setminus O$.

We apply the RoVista rule of \Cref{sec:prop}, with the latest measurement up to the return identifier.
Of measured AS--link observations, 36--43\,\% are RoVista-protected in $R\setminus O$, against 13--18\,\% in $R\cap O$.
RoVista covers 61--77\,\% of them.

The share in $R\setminus O$ is in the same range as among the carriers of ordinary valid routes in the frontier of \Cref{sec:prop}, 0.33--0.41.
The low share is thus in $R\cap O$, the ASes that carried the invalid route.
This fits ROV, but carriage and RoVista protection both measure reaching invalid prefixes.
It does not show that a given AS filtered the route itself~\cite{li2026rscope}.
% src: docs/paper/corrected-recovery-2026-10-16.json years.*.overlap, rov; generated/recovery_pairs.tex
